\title{Controlling Text-to-Image Diffusion by Orthogonal Finetuning}

\begin{abstract}
State-of-the-art text-to-image diffusion models demonstrate remarkable proficiency in synthesizing photorealistic visuals from natural language descriptions. A significant open research question is how to systematically steer or modulate these models to accomplish a variety of specialized tasks. In this context, we present Orthogonal Finetuning~(OFT), a methodologically grounded approach to adapting text-to-image diffusion models for diverse downstream applications. Distinct from prior finetuning strategies, OFT guarantees the preservation of hyperspherical energy, an attribute encapsulating the inter-neuronal relationships on the unit hypersphere. Our analyses reveal that maintaining this property is essential for upholding the semantic fidelity in generative outputs of such models. To bolster the robustness of the finetuning process, we introduce Constrained Orthogonal Finetuning (COFT), which adds a supplementary radius constraint to the hypersphere. In particular, our investigations focus on two pivotal text-to-image finetuning tasks: subject-driven generation—aiming to create images of a specific subject, based on a small set of examples and a descriptive prompt; and controllable generation—enabling models to interpret auxiliary control signals. Experimental results indicate that the OFT paradigm surpasses existing techniques in both output quality and convergence efficiency.
\end{abstract}

\section{Introduction}

Modern text-to-image diffusion models~\cite{saharia2022photorealistic, ramesh2022hierarchical, rombach2022high} deliver outstanding results in synthesizing high-fidelity images under textual direction. Nonetheless, such text conditioning often lacks the precision necessary for detailed and nuanced control over the synthesized images. To bridge this gap, our work focuses on the following two primary categories of text-to-image generation tasks:

\begin{itemize}[leftmargin=*,nosep]

    \item \textbf{Subject-driven generation}~\cite{ruiz2023dreambooth}: Here, given only a limited set of images depicting a particular subject, the goal is to produce new renditions of that subject in varied contexts through textual prompts.
    \item \textbf{Controllable generation}~\cite{zhang2023adding, mou2023t2i}: In this task, supplementary control inputs (such as canny edge detections or segmentation masks) are provided alongside a text prompt, and the objective is to generate images in accordance with both the control signal and textual description.
\end{itemize}

Both tasks fundamentally concern the challenge of fine-tuning text-to-image diffusion models in a manner that retains the generative capabilities acquired during pretraining. The criteria for a successful fine-tuning strategy can be summarized as: (1) \emph{training efficiency}—minimizing the number of parameters that need adjustment and reducing the epochs required for training; and (2) \emph{generalizability preservation}—maintaining high-quality, diverse generative outputs. Fine-tuning approaches typically involve either direct adjustment of neuron weights using a modest learning rate (\emph{e.g.}, \cite{ruiz2023dreambooth}), or the addition of compact modules with reparameterized weights (\emph{e.g.}, \cite{hulora2022,zhang2023adding}). Nevertheless, neither approach can robustly ensure that generative performance from pretraining remains intact. Furthermore, the field lacks a solid theoretical foundation for formulating effective fine-tuning schemes and for choosing optimal hyperparameters such as training duration. This challenge is compounded by the absence of robust metrics for assessing how much of the original generative ability is preserved post-finetuning. In practice, current techniques operate under the implicit assumption that reducing the Euclidean distance between the fine-tuned model and its pretrained predecessor correlates with better retention of pretrained functions. As a consequence, fine-tuning is often conducted with minimal learning rates. While this hypothesis may be valid in some cases, we contend that merely considering the Euclidean distance to the pretrained weights does not adequately represent semantic consistency. Thus, adopting a more structurally informed metric to quantify differences between the fine-tuned and pretrained models may yield significant improvements both in preserving the pretraining capabilities and enhancing the stability of fine-tuning procedures.

Building on the empirical findings that hyperspherical similarity effectively captures semantic relationships~\cite{liu2017deep,liu2018decoupled,chen2020angular}, we utilize hyperspherical energy~\cite{liu2018learning} as a metric to represent the pairwise relationships among neurons within a layer. Here, hyperspherical energy is quantified as the aggregate of hyperspherical similarities (\emph{e.g.}, cosine similarity) between all neuron pairs, reflecting how uniformly distributed the neurons are over the unit hypersphere~\cite{liu2021learning}. Our working hypothesis is that an optimally finetuned model should maintain hyperspherical energy close to that of the pretrained counterpart. One simple strategy would be to employ a regularization term that enforces constant hyperspherical energy throughout finetuning; however, this approach does not necessarily ensure that the energy discrepancy is minimized effectively. To overcome this limitation, we leverage a key invariance: hyperspherical energy remains unchanged under uniform orthogonal transformations applied to all neurons. Based on this invariance property, we introduce Orthogonal Finetuning~(OFT), a method designed to adapt large-scale text-to-image diffusion models to new tasks while preserving their hyperspherical energy. The core principle of OFT is to learn a single, shared orthogonal transformation for neurons in a layer so that their pairwise angular relationships are maintained. This process can also be interpreted as rotating the canonical coordinate system shared by the neurons in a given layer. Furthermore, taking into account that maintaining a small Euclidean distance between finetuned and pretrained models tends to better retain the advantages of pretraining, we propose a refined variant, Constrained Orthogonal Finetuning~(COFT). COFT imposes an additional constraint, restricting the finetuned model to remain within a fixed-radius hypersphere centered at the pretrained neuron configurations.

The rationale behind utilizing orthogonal transformation for neuron finetuning is partially inspired by the principles of 2D Fourier transforms, which enable the decomposition of an image into its magnitude and phase spectra. Importantly, the phase spectrum encapsulates the angular relations between the input and the basis functions, thereby preserving the majority of semantic content. Illustratively, one can reconstruct an image retaining its semantic essence by combining its phase spectrum with an arbitrary magnitude spectrum. This observation implies that altering the orientation of neurons is fundamental to achieving semantic modifications in generated images—an objective central to both subject-driven and controllable generative tasks. Nevertheless, unrestricted changes in neuron orientations risk undermining the generative capabilities acquired during pretraining. To address this, we introduce a constraint that maintains the angles between all pairs of neurons, guided by the assumption that such angular relationships underpin the neural network’s internal representation of knowledge. This perspective naturally motivates the adoption of a layer-wise orthogonal transformation that is shared within each layer, thereby keeping the hyperspherical energy invariant.

Additionally, we are influenced by work on orthogonal over-parameterized training~\cite{liu2021orthogonal}, which initializes classification neural networks randomly and employs orthogonal transformations throughout training. This methodology leverages the theoretical result that a randomly initialized network is expected to exhibit low hyperspherical energy, a quantity that \cite{liu2021orthogonal} aims to maintain throughout training—since low energy has been shown to correlate with improved generalization in classification tasks~\cite{liu2018learning,lin2020regularizing}. The findings in \cite{liu2021orthogonal} establish that orthogonal transformations possess sufficient expressivity to train neural networks with robust generalization capabilities. By contrast, our focus lies in finetuning text-to-image diffusion models to achieve more precise control and enhanced performance in downstream generation tasks. We highlight two primary distinctions between OFT and the method proposed by \cite{liu2021orthogonal}. First, while \cite{liu2021orthogonal} targets minimization of hyperspherical energy, OFT’s objective is to preserve the hyperspherical energy inherited from the pretrained model, thereby safeguarding the model’s intrinsic structure from degradation during finetuning. Particularly in the context of diffusion models, an aggressive reduction of hyperspherical energy may be detrimental to original semantic information. Second, OFT is formulated to limit deviations from the pretrained configuration, resulting in a constrained approach. Conversely, \cite{liu2021orthogonal} imposes no analogous restriction. The essence of effective finetuning lies in balancing adaptability and preservation of model structure; we posit that the OFT methodology strikes this balance adeptly. Our key contributions are as follows:

\begin{itemize}[leftmargin=*,nosep]

    \item We introduce Orthogonal Finetuning (OFT), a new approach for refining text-to-image diffusion models to enhance controllability. To address concerns regarding stability, we further design a constrained variant of OFT that restricts the model’s angular shift relative to its pretrained parameters.
    \item Unlike existing finetuning strategies, OFT is unique in maintaining the hyperspherical energy during model updates, which our empirical analysis identifies as a critical factor for preserving the pretrained model’s generative semantics.
    \item We deploy OFT on two applications: subject-driven generation and controllable generation. Our extensive experimental evaluation reveals that OFT consistently outperforms prior techniques, exhibiting superior image quality, faster convergence, and greater finetuning stability. Furthermore, OFT demonstrates improved data efficiency by achieving strong results with significantly fewer training samples and epochs.
    \item For evaluating controllability, we propose a novel control consistency metric. This metric measures the alignment between the control signal inferred from the output image and the original control directive, providing quantitative evidence of OFT's high degree of controllability.
\end{itemize}

\section{Related Work}

\textbf{Text-to-image diffusion models}. Recent years have seen remarkable advances~\cite{saharia2022photorealistic,ramesh2022hierarchical,rombach2022high,gu2022vector,nichol2022glide} in the area of generating images from text, primarily fueled by the progress in diffusion-based generative modeling~\cite{ho2020denoising,song2021score,song2021denoising,dhariwal2021diffusion} and advances in learning multimodal visual-language representations~\cite{su2020vlbert,lu2019vilbert,radford2021learning,wang2021simvlm,li2022blip,li2023blip,alayrac2022flamingo,schuhmann2022laion}. GLIDE~\cite{nichol2022glide} and Imagen~\cite{saharia2022photorealistic} build diffusion models directly in pixel space; the former simultaneously optimizes a text encoder alongside a diffusion model on paired data, whereas the latter relies on a pretrained, frozen text encoder. In contrast, Stable Diffusion~\cite{rombach2022high} and DALL-E2~\cite{ramesh2022hierarchical} operate in a learned latent space. Stable Diffusion leverages VQ-GAN~\cite{esser2021taming} to construct a visual codebook for its latent representations, and DALL-E2 uses CLIP~\cite{radford2021learning} to develop a unified latent embedding space for both modalities. Beyond diffusion models, generative adversarial networks~\cite{reed2016generative,zhang2017stackgan,xu2018attngan,li2019controllable} and autoregressive models~\cite{ramesh2021zero,ding2021cogview,wu2022nuwa,yuscaling} have also contributed to progress in this field. OFT’s design as a model-agnostic finetuning technique enables its integration with any text-to-image diffusion model architecture.

\textbf{Subject-driven generation}. To mitigate unwanted alterations of the main subject, several works~\cite{avrahami2022blended,nichol2022glide} incorporate user-specified masks as auxiliary conditions. Alternatively, inversion-based strategies~\cite{choi2021ilvr,dhariwal2021diffusion,ramesh2022hierarchical,gal2022image} allow for contextual changes without disturbing the identity of the subject. The approach by~\cite{hertz2022prompt} is capable of both local and global edits, independent of any mask input. However, most of these methods fall short in preserving fine-grained identity-specific features. Pivotal Tuning~\cite{roich2022pivotal} refines the generator around an initially inverted code and introduces an extra regularization term to maintain subject identity. Along the same lines,~\cite{nitzan2022mystyle} constructs a personalized generative prior from a dataset of facial images belonging to one individual. In contrast,~\cite{casanova2021instance} enables diverse instance-level generation, though instance-specific details may be compromised. DreamBooth~\cite{ruiz2023dreambooth} fine-tunes a text-to-image diffusion model using both a reconstruction loss and a class-specific prior, employing a dedicated token and a handful of subject images. While OFT is built upon the DreamBooth setup, it replaces standard low-rate finetuning with updates constrained by orthogonal transformations.

\textbf{Controllable generation}. Image-to-image translation is essentially a subset of controllable generation tasks. Earlier works in this area largely employ conditional generative adversarial networks~\cite{isola2017image,zhu2017toward,wang2018high,park2019semantic,choi2018stargan}. In addition, diffusion models have been utilized for this translation problem~\cite{saharia2022palette,voynov2022sketch,wang2022pretraining}. More recent developments, such as ControlNet~\cite{zhang2023adding}, enhance controllability by finetuning existing diffusion models with supplementary control signals, thereby achieving highly effective and controllable generation. Similarly, T2I-Adapter~\cite{mou2023t2i}, introduced independently, also targets increased control in diffusion-based image generation by finetuning pretrained models. Adopting the experimental setup described in \cite{zhang2023adding,mou2023t2i}, our approach leverages OFT for finetuning diffusion models. This consistently results in enhanced controllability, even when using smaller amounts of training data and fewer finetuning parameters. Importantly, OFT adds no extra computational cost during test-time inference.

\textbf{Model finetuning}. The strategy of adapting large pretrained models to specific tasks via finetuning has become widespread~\cite{devlin2018bert,brown2020language,he2022masked}. Various adaptation techniques (\emph{e.g.}, \cite{hulora2022,houlsby2019parameter,pfeiffer2020adapterfusion}) have been extensively explored in natural language processing. Of particular relevance is LoRA~\cite{hulora2022}, which presumes that additive updates to the weights during finetuning are of low rank. Conversely, OFT employs a multiplicative, layer-shared orthogonal transformation for updating neuron weights. This technique mathematically guarantees the preservation of pairwise angular relationships among neurons within a layer, leading to improved stability.

\section{Orthogonal Finetuning}

\subsection{Why Does Orthogonal Transformation Make Sense?}

To motivate the use of orthogonal transformation in the context of finetuning text-to-image diffusion models, we break the issue down into two key points: (1) the motivation for adjusting the angles (directions) of neurons during finetuning, and (2) the justification for utilizing orthogonal transformations as the mechanism to update these angles.

To address the first question, we take cues from the findings in \cite{liu2018decoupled,chen2020angular}, where it is shown that differences in angular features effectively capture the semantic gap. SphereNet~\cite{liu2017deep} demonstrates that constraining all neurons of a neural network to lie on a unit hypersphere preserves the model’s representational power and even enhances generalization, suggesting that the orientation of neurons is sufficient for encoding salient information from data. To further underscore the significance of neuron orientation, we present a simplified experiment in Figure~\ref{fig:toy_ae} in which a basic convolutional autoencoder is trained on flower images. During training, standard inner product computation is used to generate the feature map ($z$ is the output corresponding to the convolution kernel $\bm{w}$ applied to input $\bm{x}$ in each local patch). For evaluation, we test three alternative methods of generating the feature map: (a) using the inner product as during training, (b) isolating magnitude, and (c) isolating angular components. The visual results, displayed in Figure~\ref{fig:toy_ae}, indicate that the angular component of neuron activations can nearly completely reconstruct the original input, while the magnitude component lacks informative content. It is notable that we did not use the cosine-based (angular) computation (c) during training; the training process relied exclusively on inner products. These findings point to the dominant role of neuron directions (angles) in representing the semantic content of images, indicating that adjusting neuron orientations may be especially effective for altering semantic characteristics.

Regarding the second question, a direct approach to finetuning neuron directions is to jointly rotate or reflect all neurons within a layer, a process inherently linked to orthogonal transformations. While more sophisticated angular transformations—such as rotating each neuron by a distinct angle—might offer additional flexibility, we have observed that orthogonal transformations achieve a desirable balance between adaptability and structure. Further, prior work \cite{liu2021orthogonal} has established the effectiveness of orthogonal transformations in learning for neural networks. To substantiate this perspective, we carry out an experiment examining the regularization provided by the orthogonality constraint using the setup from ControlNet~\cite{zhang2023adding}. The corresponding outcomes, displayed in Figure~\ref{fig:ortho}, reveal that our standard Orthogonal Fine-Tuning (OFT) consistently achieves robust and precise controllability after completion of training (at epoch 20). In contrast, an OFT variant lacking the orthogonality restriction is unable to generate realistic images or provide effective control. This experimental evidence highlights the essential role of orthogonality constraints in the efficacy of OFT.

\subsection{General Framework}

Finetuning is conventionally performed via gradient descent with a small step size, often only adjusting specific model layers. Utilizing a small learning rate serves to constrain updates, maintaining close proximity to the pretrained parameters; standard finetuning thereby minimizes $\thickmuskip=2mu \medmuskip=2mu \| \bm{M} - \bm{M}^0\|$ where $\bm{M}$ and $\bm{M}^0$ denote the weights after and before finetuning, respectively. While this implicit regularization aligns with the intuition of preserving the pretrained model's learned features, it may still allow excessive freedom, particularly for large-scale models. To further restrict the solution space, LoRA applies a low-rank constraint to parameter adjustments, imposing $\thickmuskip=2mu \medmuskip=2mu \text{rank}(\bm{M} - \bm{M}^0)=r'$ where $r'$ is a predefined small integer. In contrast, OFT takes a different approach by constraining the hyperspherical energy difference between neurons, enforcing $\thickmuskip=2mu \medmuskip=2mu \| \text{HE}(\bm{M})-\text{HE}(\bm{M}^0) \|=0$, where $\text{HE}(\cdot)$ calculates the hyperspherical energy for a set of weights.

To illustrate, consider a fully connected layer $\thickmuskip=2mu \medmuskip=2mu \bm{W}=\{\bm{w}_1,\cdots,\bm{w}_n\}\in\mathbb{R}^{d\times n}$ with columns $\bm{w}_i\in\mathbb{R}^{d}$ representing the $i$-th neuron's weights (and $\bm{W}^0$ as the initialization). The output $\thickmuskip=2mu \medmuskip=2mu \bm{z}\in\mathbb{R}^n$ is computed as $\thickmuskip=2mu \medmuskip=2mu \bm{z}=\bm{W}^\top\bm{x}$ for an input $\thickmuskip=2mu \medmuskip=2mu \bm{x}\in\mathbb{R}^d$. Within this framework, OFT may be interpreted as seeking to minimize the gap in hyperspherical energy between the adapted and original weights:

\begin{equation}\label{eq:oft_principle}
\footnotesize
\min_{\bm{W}} \left\| \text{HE}(\bm{W})-\text{HE}(\bm{W}^0)\right\|~~\Leftrightarrow~~\min_{\bm{W}} \bigg{\|} \sum_{i\neq j} \|\hat{\bm{w}}_i-\hat{\bm{w}}_j\|^{-1}-\sum_{i\neq j} \|\hat{\bm{w}}^0_i-\hat{\bm{w}}^0_j\|^{-1}\bigg{\|}
\end{equation}

Here, each normalized neuron is given by $\thickmuskip=2mu \medmuskip=2mu \hat{\bm{w}}_i:=\bm{w}_i/\|\bm{w}_i\|$, and hyperspherical energy is $\thickmuskip=2mu \medmuskip=2mu \text{HE}(\bm{W}):= \sum_{i\neq j} \|\hat{\bm{w}}_i-\hat{\bm{w}}_j\|^{-1}$. Notably, the minimum value for Eq.~\eqref{eq:oft_principle} is attainable when $\bm{W}$ and $\bm{W}^0$ are related via a rigid transformation—specifically, if $\thickmuskip=2mu \medmuskip=2mu \bm{W}=\bm{R}\bm{W}^0$, with $\thickmuskip=2mu \medmuskip=2mu \bm{R}\in\mathbb{R}^{d\times d}$ being an orthogonal matrix (where the sign of the determinant differentiates between rotation and reflection). This principle underpins OFT: model finetuning is accomplished solely by learning an orthogonal transformation shared across neurons in each layer. More precisely, for a given pretrained fully connected layer $\thickmuskip=2mu \medmuskip=2mu\bm{W}^0\in

\begin{equation}\label{eq:oft_general}
    \footnotesize
    \bm{z} = \bm{W}^\top\bm{x} = (\bm{R}\cdot\bm{W}^0)^\top\bm{x},~~~\text{s.t.}~\bm{R}^\top\bm{R} = \bm{R}\bm{R}^\top = \bm{I}
\end{equation}

Here, $\bm{W}$ refers to the weight matrix obtained after OFT fine-tuning, and $\bm{I}$ represents the identity matrix. Figure~\ref{fig:block} visually outlines the OFT framework. Analogous to the initialization with zeros as employed in LoRA, it is essential for OFT to enable fine-tuning to commence precisely from the original pretrained parameters $\bm{W}^0$. This is accomplished by setting the orthogonal matrix $\bm{R}$ initially as the identity matrix, thereby ensuring the finetuned parameters coincide with the pretrained ones at the outset. To preserve the orthogonal property of $\bm{R}$, various differentiable orthogonalization approaches are available, such as those proposed by \cite{liu2021orthogonal,lezcano2019cheap}. In subsequent sections, we address strategies for maintaining orthogonality in a way that is computationally practical.

\subsection{Efficient Orthogonal Parameterization}\label{sect:eff_ortho}

Traditional orthogonalization techniques, like the Gram-Schmidt procedure, while differentiable, tend to be computationally intensive in practical applications~\cite{liu2021orthogonal}. To improve computational efficiency, we employ Cayley parameterization to generate the required orthogonal matrices. More specifically, the orthogonal matrix is expressed as $\thickmuskip=2mu \medmuskip=2mu \bm{R}=(\bm{I}+\bm{Q})(I-\bm{Q})^{-1}$, where the matrix $\bm{Q}$ is skew-symmetric, fulfilling $\thickmuskip=2mu \medmuskip=2mu \bm{Q} = -\bm{Q}^\top$. This method enhances efficiency but restricts the result to orthogonal matrices of determinant $1$, placing them in the special orthogonal group. Nonetheless, our experiments indicate that this constraint does not impede empirical performance. Even with the Cayley parameterization, constructing a fully dense orthogonal matrix $\bm{R}$ is highly parameter-intensive when $d$ is large. To alleviate this, we advocate using a block-diagonal structure for $\bm{R}$ composed of $r$ blocks, resulting in:

\begin{equation}
    \footnotesize
    \bm{R} = \text{diag}(\bm{R}_1, \bm{R}_2, \cdots, \bm{R}_r) = \begin{bmatrix}
    \bm{R}_1 \in \text{O}(\frac{d}{r}) &  &\\
    &\ddots & \\
    & & \bm{R}_r \in \text{O}(\frac{d}{r})
    \end{bmatrix} \in \text{O}(d)
\end{equation}

Here, $\text{O}(d)$ refers to the group of orthogonal matrices in $d$ dimensions, where $\thickmuskip=2mu \medmuskip=2mu \bm{R}\in\mathbb{R}^{d\times d}$ and for all $i$, $\thickmuskip=2mu \medmuskip=2mu \bm{R}_i\in\mathbb{R}^{d/r\times d/r}$ are orthogonal matrices. When the parameter $r$ is set to $1$, the block-diagonal structure simplifies, making the orthogonal matrix fully unconstrained. In the case of an orthogonal matrix of size $\thickmuskip=2mu \medmuskip=2mu d\times d$, the total number of independent parameters is $\thickmuskip=2mu \medmuskip=2mu d(d-1)/2$, leading to a computational complexity of $\mathcal{O}(d^2)$. For a block-diagonal orthogonal matrix partitioned into $r$ blocks, the number of parameters becomes $\thickmuskip=2mu \medmuskip=2mu d(d/r-1)/2$, resulting in $\mathcal{O}(d^2/r)$ complexity. Moreover, by enforcing that all block matrices are shared, i.e., $\thickmuskip=2mu \medmuskip=2mu\bm{R}_i=\bm{R}_j$ for every pair $i\neq j$, the parameter count can be reduced further, now scaling as $\mathcal{O}(d^2/r^2)$. Importantly, all such approaches continue to guarantee that $\bm{R}$ remains an orthogonal matrix, and thus, the hyperspherical structure is exactly maintained.

We next compare the parameter economy of OFT with that of LoRA. For a LoRA module characterized by a rank parameter $r'$, the number of its trainable weights is given by $\thickmuskip=2mu \medmuskip=2mu r'(d+n)$. If we tie both $r$ and $r'$ linearly to $d$ (for example, setting $\thickmuskip=2mu \medmuskip=2mu r=r'=\alpha d$, with $\thickmuskip=2mu \medmuskip=2mu 0<\alpha\leq 1$), the parameter complexity for LoRA evaluates to $\mathcal{O}(d^2+dn)$, whereas for OFT it drops to $\mathcal{O}(d)$. For illustration, take a $128\times 128$ weight matrix: employing LoRA with $\thickmuskip=2mu \medmuskip=2mu r'=8$ results in $2,048$ tunable parameters, while OFT using $\thickmuskip=2mu \medmuskip=2mu r=8$ (without block sharing) involves only $960$ trainable parameters.

\subsection{Constrained Orthogonal Finetuning}

The expressiveness of OFT can be further reduced by restricting the optimization so that the updated model does not deviate much from the original pretrained parameters. In this setup, COFT adopts the following computation:

\begin{equation}\label{eq:coft_general}
    \footnotesize
    \bm{z}=\bm{W}^\top\bm{x}=(\bm{R}\cdot\bm{W}^0)^\top\bm{x},~~~\text{s.t.}~\bm{R}^\top\bm{R}=\bm{R}\bm{R}^\top=\bm{I},~\norm{\bm{R}-\bm{I}}\leq \epsilon
\end{equation}

Here, the formulation enforces both that $\bm{R}$ is orthogonal and that it remains close (within $\epsilon$) to the identity matrix. While the orthogonality can be enforced using the Cayley parameterization as described in Section~\ref{sect:eff_ortho}, imposing the $\epsilon$-norm restriction is challenging when the Cayley transform is utilized. To provide intuition into this issue, recall the Cayley transform expresses $\thickmuskip=2mu \medmuskip=2mu \bm{R}=(\bm{I}+\bm{Q})(I-\bm{Q})^{-1}$. When approximated using the Neumann series expansion, this yields $\thickmuskip=2mu \medmuskip=2mu \bm{R}\approx\bm{I}+2\

Because the series converges with respect to the operator norm, the constraint $\thickmuskip=2mu \medmuskip=2mu\|\bm{R}-\bm{I}\|\leq\epsilon$ can be reformulated inside the Cayley transform. Under this transformation, it becomes an equivalent restriction $\thickmuskip=2mu \medmuskip=2mu \|\bm{Q}-\bm{0}\|\leq\epsilon'$, where $\bm{0}$ is the zero matrix and $\epsilon'$ serves as a distinct hyperparameter controlling error (not the same as $\epsilon$). Enforcing this new constraint on $\bm{Q}$ is straightforward using projected gradient descent methods. To guarantee that the orthogonal matrix $\bm{R}$ starts as the identity, we initialize $\bm{Q}$ to be the zero matrix. Conceptually, COFT integrates two explicit regularizations: minimizing hyperspherical energy divergence and imposing a bound on the deviation from the pretrained parameters. The latter constraint is commonly incorporated in standard finetuning techniques in an implicit manner, while COFT incorporates it explicitly. Our results show that although OFT achieves strong performance, the added explicit deviation constraint in COFT leads to even more robust finetuning outcomes compared to OFT alone. Figure~\ref{fig:eps_coft} illustrates the influence of $\epsilon$ on COFT's effectiveness. The figure reveals that $\epsilon$ modulates the degree of adaptation allowed during finetuning: with a higher $\epsilon$, COFT produces results similar to OFT; with a lower $\epsilon$, the finetuned model retains more characteristics of the original pretrained text-to-image diffusion model.

\subsection{Re-scaled Orthogonal Finetuning}

We introduce a straightforward augmentation to standard OFT: the inclusion of a neuron-wise scaling parameter. This enhancement is motivated by the observation that scaling neuron activations does not alter their hyperspherical energy, since normalization removes magnitude information. Our modified forward computation is: $\thickmuskip=2mu \medmuskip=2mu\bm{z}=(\bm{R}\bm{W}^0\bm{D})^\top\bm{x}$\footnote{\textbf{Errata}: The forward computation in the NeurIPS camera-ready submission was incorrectly specified as $\thickmuskip=2mu \medmuskip=2mu\bm{z}=(\bm{D}\bm{R}\bm{W}^0)^\top\bm{x}$. The actual code implementation is correct, hence the results in Appendix~\ref{app:rescaled} are unaffected.}, where the scaling matrix $\thickmuskip=2mu \medmuskip=2mu \bm{D}=\text{diag}(s_1,\cdots,s_n)\in\mathbb{R}^{n\times n}$ is diagonal and learnable, with each diagonal entry $s_1,\ldots,s_n$ positive. Unlike the original OFT forward pass in Eq.~\eqref{eq:oft_general}, which optimizes only $\bm{R}$, both $\bm{D}$ and $\bm{R}$ are now trainable. This re-scaled version enhances OFT's expressiveness with a negligible increase in parameter count. Nonetheless, our experiments focus on the vanilla OFT to isolate the impact of orthogonal transformation. We observe, however, that re-scaled OFT generally performs better; details are provided in Appendix~\ref{app:rescaled}.

\section{Intriguing Insights and Discussions}

\textbf{OFT’s architecture independence}. OFT is fundamentally compatible with any neural network architecture. In the Transformer context, LoRA is generally deployed on attention weight matrices~\cite{hulora2022}, so for direct comparison, we restrict OFT finetuning to the same attention weights in our experiments. Furthermore, OFT extends naturally to convolutional layers. Its block-diagonal matrix $\bm{R}$ enables intuitive applications in convolution settings (unlike LoRA). For example, when the number of blocks matches the number of input channels, each block individually modifies a single neuron channel—akin to depth-wise convolution kernel training~\cite{chollet2017xception}. When all blocks in $\bm{R}$ are tied, the same orthogonal transformation is applied across channels. Preliminary exploration of OFT for convolutional layer finetuning is provided in Appendix~\ref{app:convolution}.

\textbf{Connection to LoRA}. LoRA stabilizes the information contained in pretrained weight matrices by introducing a low-rank update, thereby limiting major deviations during finetuning. Conversely, OFT enforces an orthogonal (full-rank) transformation of the pretrained weights, which serves to protect the integrity of the original learned knowledge by preserving orthogonality. The OFT forward pass can be re-expressed as $\thickmuskip=2mu \medmuskip=2mu \bm{z} = (\bm{R}\bm{W}^0)^\top \bm{x} = (\bm{W}^0 + (\bm{R} - \bm{I})\bm{W}^0)^\top \bm{x}$, where the term $\thickmuskip=2mu \medmuskip=2mu (\bm{R} - \bm{I})\bm{W}^0$ functions analogously to LoRA's low-rank modification. Since $\bm{W}^0$ is typically full-rank, OFT can also achieve a low-rank adaptation if $\thickmuskip=2mu \medmuskip=2mu \bm{R} - \bm{I}$ is chosen to be low-rank. Parallel to LoRA’s rank parameter $r'$, OFT utilizes a block-diagonal size parameter $r$ to manage the number of trainable parameters. Interestingly, the two approaches embody fundamentally different parameter-efficient strategies: while LoRA utilizes low-rank decomposition to economize parameter count, OFT leverages block-wise sparsity through orthogonality.

\textbf{Why OFT converges faster}. As illustrated in Figure~\ref{fig:toy_ae}, altering the orientation of neurons constitutes the most efficient means for semantic modification, which precisely aligns with the design motivation of OFT. Additionally, OFT may be interpreted as performing finetuning constrained to the surface of a smooth hypersphere, potentially offering a more favorable optimization landscape. Empirical findings in \cite{liu2021orthogonal} corroborate this hypothesis.

\textbf{Why not minimize hyperspherical energy}. A notable distinction from \cite{liu2021orthogonal} is that our objective does not include hyperspherical energy minimization. In classification contexts, it is desirable to avoid redundant neurons. Minimizing hyperspherical energy results in neurons being distributed uniformly over the hypersphere, which is not ideal for finetuning, as such uniformity may disrupt the beneficial structure acquired during pretraining.

\textbf{Balancing flexibility and regularization in finetuning}. Our investigation reveals a fundamental tension between flexibility and regularization. While conventional finetuning offers maximal flexibility, it often results in instability and is susceptible to model collapse. OFT, notable for its simplicity, effectively navigates this trade-off by maintaining the pairwise angles between neurons, thereby achieving an appropriate compromise. In contrast, the block-diagonal parameterization serves as a more stringent form of regularization on the orthogonal matrix.

\textbf{No increase in inference computation}. Distinct from approaches like ControlNet, OFT does not impose extra computational cost at inference. During the inference phase, the optimized orthogonal matrix $\bm{R}$ is directly multiplied by the pretrained weights $\bm{W}^0$, forming a new effective parameter matrix $\thickmuskip=2mu \medmuskip=2mu \bm{W}=\bm{R}\bm{W}^0$. As a result, inference operates with the same efficiency as the original pretrained network.

\section{Experiments and Results}

\textbf{Experimental configuration}. Our experiments are built on the Stable Diffusion v1.5~\cite{rombach2022high} architecture for text-to-image synthesis. To ensure equitable evaluation, we sample generated images from all methods randomly. The procedures for subject-driven generation mirror those from DreamBooth~\cite{ruiz2023dreambooth}, while our setup for controllable generation is aligned with ControlNet~\cite{zhang2023adding} and T2I-Adapter~\cite{mou2023t2i}. To fairly benchmark OFT and COFT against LoRA, modifications are confined to the same network layers where LoRA operates. Additional experimental specifics and comprehensive results can be found in Appendix~\ref{app:settings}.

\subsection{Subject-driven Generation}

\textbf{Experimental setup}. The baselines include DreamBooth~\cite{ruiz2023dreambooth} and LoRA~\cite{hulora2022}. All approaches utilize the loss function specified in DreamBooth. For DreamBooth and LoRA, our methodology adheres to the original protocols, with hyperparameters chosen based on recommended optimal configurations. Further results and implementation details are reported in Appendix~\ref{app:settings},\ref{app:coft_oft},\ref{app:more_qualitative},\ref{app:failure}.

\textbf{Finetuning stability and convergence}. To begin, we assess the stability and speed of convergence during finetuning for DreamBooth, LoRA, OFT, and COFT. The outcomes are presented in Figure~\ref{fig:teaser} and Figure~\ref{fig:db_convergence}. It is evident that both COFT and OFT maintain stable finetuning processes for the diffusion model. At 400 training steps, effective control is demonstrated by DreamBooth and both OFT variants, whereas LoRA struggles to retain subject identity. As training proceeds to 2000 iterations, DreamBooth encounters mode collapse and starts generating defective images; LoRA is unable to synthesize a yellow shirt and instead produces yellow fur. Conversely, OFT and COFT maintain both stable and reliable control over the generated outputs. These findings support the rapid yet robust convergence properties of the OFT framework in scenarios involving subject-driven image generation. Importantly, insensitivity to the total number of finetuning steps simplifies practical application, as it lessens the need for meticulous adjustment per subject. With both OFT and COFT, a larger fixed iteration number can be applied broadly without additional hyperparameter tuning. For COFT, selecting an appropriate $\epsilon$ makes setting both the learning rate and iteration number straightforward.

\textbf{Quantitative comparison}. In line with the methodology of \cite{ruiz2023dreambooth}, we perform a quantitative evaluation of subject fidelity (using DINO~\cite{caron2021emerging} and CLIP-I~\cite{radford2021learning}), text prompt fidelity (using CLIP-T~\cite{radford2021learning}), and sample diversity (with LPIPS~\cite{zhang2018unreasonable}). The CLIP-I metric measures the mean pairwise cosine similarity between CLIP embeddings for generated and real images. DINO provides a comparable assessment, differing only in that it utilizes ViT S/16 DINO embeddings. CLIP-T computes the mean cosine similarity between CLIP embeddings of the generated images and their corresponding text prompts. Sample diversity is quantified by the average LPIPS cosine similarity among images generated for the same subject given an identical text prompt. As seen in Table~\ref{table:dreambooth_final}, COFT and OFT substantially outperform both DreamBooth and LoRA in DINO and CLIP-I scores, while their results for text prompt fidelity and diversity metrics are slightly superior or similar. Each approach is evaluated by repeatedly finetuning a common pretrained model using 30 distinct random seeds to control for variability. Overall, the data indicate that our OFT framework consistently yields superior convergence, robustness, and final performance compared to existing baselines.

\textbf{Qualitative comparison}. To provide a clearer illustration of OFT's advantages, we present a selection of randomly chosen cases demonstrating subject-driven generation in Figure~\ref{fig:dreambooth_final}. For equitable evaluation, all outputs are produced by the same finetuned model with different methods, ensuring that text prompts are not individually optimized for specific outputs. For each approach, we report the results from the model exhibiting the highest validation CLIP scores. Inspection of Figure~\ref{fig:dreambooth_final} indicates that both OFT and COFT maintain strong semantic fidelity to the subject, whereas LoRA frequently fails to retain subject identity (for instance, the LoRA output in the bowl scenario loses the subject altogether). Furthermore, both OFT and COFT enable more precise manipulation via textual inputs, in contrast to DreamBooth, which, despite effectively preserving subject identity, frequently struggles to generate outputs aligned with the provided prompt (as seen in the top row of the bowl example). This qualitative assessment highlights OFT’s ability to concurrently realize enhanced controllability and subject retention. Additionally, OFT demonstrates robustness with respect to iteration count, maintaining high performance even as the number of iterations increases, a property not observed in DreamBooth or LoRA. Additional qualitative illustrations are presented in Appendix~\ref{app:more_qualitative}. Human evaluation results in Appendix~\ref{app:human} provide further evidence in support of OFT's superior performance.

\subsection{Controllable Generation}

\textbf{Settings}. For baselines, we consider ControlNet~\cite{zhang2023adding}, T2I-Adapter~\cite{mou2023t2i}, and LoRA~\cite{hulora2022}. The main paper investigates three complex controllable generation scenarios: Canny edge to image~(C2I) using the COCO dataset~\cite{lin2014microsoft}, segmentation map to image~(S2I) using the ADE20K dataset~\cite{zhou2017scene}, and landmark to face~(L2F) using the CelebA-HQ dataset~\cite{karras2018progressive,xia2021tedigan}. Each method is employed to finetune Stable Diffusion~(SD) v1.5 across these datasets for 20 training epochs. Extended results can be found in Appendix~\ref{app:more_qualitative},~\ref{app:more_control_task}, and~\ref{app:failure}.

\textbf{Convergence}. We assess the rate of convergence for ControlNet, T2I-Adapter, LoRA, and COFT in the context of the L2F task, employing both quantitative and qualitative measures. For the quantitative metric, we calculate the average $\ell_2$ distance between the control and predicted face landmarks. Figure~\ref{fig:face_convergence} presents the progression of face landmark error as models are fine-tuned across varying epochs. From these results, it is evident that COFT and OFT exhibit substantially accelerated convergence rates compared to other methods. For instance, LoRA requires 20 epochs to achieve an error rate comparable to that of our OFT framework at just the 8-th epoch. It is important to emphasize that both OFT and COFT maintain a trainable parameter count similar to that of LoRA—markedly lower than ControlNet—yet still deliver superior convergence efficiency relative to other techniques. The expedited convergence of OFT is further corroborated by observations in Figure~\ref{fig:teaser}, where the rightmost example demonstrates that OFT attains satisfactory convergence utilizing only 5\% of the ADE20K dataset, indicating enhanced data efficiency over ControlNet and LoRA. For qualitative analysis, we center our comparison on OFT, COFT, and ControlNet, as ControlNet’s landmark error is most comparable to ours. As shown in Figure~\ref{fig:face_convergence_qual}, both OFT and COFT display consistent and stable convergence, with generated face poses progressively aligning with the control landmarks over time. In comparison, ControlNet exhibits a sudden onset of controllability after the 8-th epoch, mirroring the abrupt convergence described in \cite{zhang2023adding}. The observed convergence stability and gradual progression of OFT not only enhance its practical usability but also yield more reliable and predictable training dynamics, thereby facilitating easier identification of optimal OFT hyperparameters.

\textbf{Quantitative comparison}. To assess the efficacy of controllable generation, we propose a control consistency metric, which involves extracting the control signal from the generated output and directly measuring its correspondence with the original input signal. Specifically, for the C2I task, we utilize IoU and F1 score as our primary evaluation metrics. The S2I task employs mean IoU as well as both mean and overall accuracy. For the L2F task, the evaluation is performed via the mean $\ell_2$ distance computed between the intended control landmarks and those predicted by the model. Further elaboration on these consistency metrics can be found in Appendix~\ref{app:settings}. Every baseline is evaluated under its optimal hyperparameter configuration. The data presented in Table~\ref{table:control_gen} demonstrates that both OFT and COFT frameworks deliver superior and more reliable control performance compared to competing techniques. Notably, we find that adapter-based methods (\emph{e.g.}, T2I-Adapter and ControlNet) not only converge more slowly but also attain inferior final outcomes. LoRA surpasses ControlNet on the S2I task but falls behind on both the C2I and L2F tasks. Overall, our investigations reveal that LoRA has a limited performance ceiling, despite exhaustive hyperparameter optimization. In contrast, OFT exhibits continual improvement as we increase the number of trainable parameters, suggesting its potential has not yet been fully realized. We would like to highlight that our approach to quantitatively evaluating controllable generation constitutes a novel contribution, enabling precise assessment of control in finetuned models—bounded only by the accuracy of the segmentation or detection models used as references.

\textbf{Qualitative comparison}. We further perform a qualitative evaluation of OFT and COFT against leading approaches such as ControlNet, T2I-Adapter, and LoRA. As shown in the randomly sampled outputs of Figure~\ref{fig:control_final}, both OFT and COFT consistently produce images with high visual realism and strong adherence to the intended control signals. Specifically, for the S2I task, LoRA is unable to synthesize images that conform to the input segmentation map, whereas ControlNet, OFT, and COFT all manage precise guidance in the image generation process. Notably, in comparison with ControlNet, OFT and COFT can deliver more detailed, photorealistic results and capture more plausible geometry, all while employing significantly fewer parameters. During the C2I task, OFT and COFT excel at generating semantically aligned images from simple Canny edge maps, outperforming T2I-Adapter and LoRA, which exhibit inferior results. Regarding the L2F task, our proposed approach achieves the highest accuracy in pose-driven face synthesis, even when handling difficult facial angles. Across these three types of control tasks, both OFT and COFT demonstrate superior qualitative performance relative to state-of-the-art counterparts, attesting to the robust controllability offered by our OFT architecture. For a broader qualitative assessment, additional results spanning all three control scenarios can be found in Appendix~\ref{app:control_exp}. We also showcase the versatility of OFT on various other control problems (such as dense pose to body rendering, sketch-to-image, and depth-to-image synthesis) in Appendix~\ref{app:more_control_task}.

\section{Concluding Remarks and Open Problems}

Building on the insight that angular relationships among neurons play a key role in shaping visual semantics, we introduce a straightforward yet potent approach—orthogonal finetuning—for regulating text-to-image diffusion models. In particular, our method is applied to two primary text-to-image generation scenarios: subject-driven generation and controllable generation. Relative to prior techniques, OFT achieves enhanced controllability and increased stability during finetuning, while utilizing a reduced number of trainable parameters. Crucially, OFT imposes no extra computational load at inference time, allowing for streamlined and practical deployment.

The development of OFT brings forth several notable research questions. First, orthogonality in OFT is maintained through Cayley parametrization, a process which necessitates matrix inversion and consequently presents a barrier to scalability. While we mitigate this limitation by employing a block diagonal parametrization, efficiently computing the matrix inverse in a manner compatible with differentiation remains unresolved. Second, OFT exhibits distinct promise for compositionality: orthogonal matrices generated from different OFT finetuning sessions can be composed via matrix multiplication, yielding another orthogonal matrix. An open question is whether these composed orthogonal matrices retain the accumulated knowledge from each individual downstream task. Lastly, the model's parameter efficiency relies heavily on the block diagonal design, a compromise that introduces certain biases and may reduce adaptability. Finding alternative strategies to enhance parameter efficiency while minimizing bias and preserving flexibility constitutes a significant open research direction.